\documentclass{styles/acmsiggraph}
\usepackage[utf8]{inputenc}
\usepackage[T1]{fontenc}
\usepackage{geometry}
\usepackage{graphicx}
\usepackage{float}
\usepackage{hyperref}
\usepackage[english,indonesian]{babel}

\title{Optimasi Pemrosesan Template LaTeX Menggunakan Pendekatan Berbasis AI}
\author{Academic Support Agent}
\date{\today}

\begin{document}
\maketitle

\begin{abstract}

\end{abstract}

\chapter{PENDAHULUAN}


\chapter{TINJAUAN PUSTAKA}


% Original content (commented out):
% \documentclass[conference]{styles/acmsiggraph}
\usepackage{comment} % enables the use of multi-line comments (\ifx \fi)
\usepackage{lipsum} %This package just generates Lorem Ipsum filler text.
\usepackage{fullpage} % changes the margin
\usepackage{enumitem} % for customizing enumerate tags
\usepackage{amsmath,amsthm,amssymb}
\usepackage{listings}
\usepackage{graphicx}
\usepackage{etoolbox}   % for booleans and much more
\usepackage{verbatim}   % for the comment environment
\usepackage[dvipsnames]{xcolor}
\usepackage{fancyvrb}
\usepackage{hyperref}
\usepackage{menukeys}
\usepackage{titlesec}
\setlength{\parskip}{.8mm}
\setcounter{MaxMatrixCols}{20}

\title{\huge Homework 6: Solutions \\ \LARGE {CSE 490V: Virtual Reality Systems}}
\author{\Large Student Name \\ student@uw.edu}
\pdfauthor{Student Name}

\hypersetup{
	colorlinks=true,
	urlcolor=[rgb]{0.97,0,0.30},
	anchorcolor={0.97,0,0.30},
	linkcolor=[rgb]{0.97,0,0.30},
	filecolor=[rgb]{0.97,0,0.30},
}

% redefine \VerbatimInput
\RecustomVerbatimCommand{\VerbatimInput}{VerbatimInput}%
{fontsize=\footnotesize,
 %
 frame=lines,  % top and bottom rule only
 framesep=2em, % separation between frame and text
 rulecolor=\color{Gray},
 %
 label=\fbox{\color{Black}\textbf{OUTPUT}},
 labelposition=topline,
 %
 commandchars=\|\(\), % escape character and argument delimiters for
                      % commands within the verbatim
 commentchar=*        % comment character
}

\titlespacing*{\section}{0pt}{5.5ex plus 1ex minus .2ex}{2ex}
\titlespacing*{\subsection}{0pt}{3ex}{2ex}

\setcounter{secnumdepth}{4}
\renewcommand\theparagraph{\thesubsubsection.\arabic{paragraph}}
\newcommand\subsubsubsection{\paragraph}

\setlength{\parskip}{0.5em}

% a macro for hiding answers
\newbool{hideanswers}
\setbool{hideanswers}{false}
\newenvironment{answer}{}{}
\ifbool{hideanswers}{\AtBeginEnvironment{answer}{\comment} %
\AtEndEnvironment{answer}{\endcomment}}{}

\newcommand{\points}[1]{\hfill \normalfont{(\textit{#1pts})}}
\newcommand{\pointsin}[1]{\normalfont{(\textit{#1pts})}}

\begin{document}
\maketitle
\section{Theoretical Part}

\subsection{Image Formation Model for Lighthouse Pose Tracking \points{10}}
\label{sec:imageformation}

Let's say we have the ground truth 6-DOF pose, i.e., orientation $\boldsymbol{\theta}$ (in degrees) and position $\vec{t}$ (in mm), of a ``VRduino'' device given as
%
$$\boldsymbol{\theta} = \begin{pmatrix}
\theta_x \\ \theta_y \\ \theta_z
\end{pmatrix} = \begin{pmatrix}
45^\circ \\ 0^\circ \\ 45^\circ
\end{pmatrix}, \quad
\vec{t} = \begin{pmatrix}
t_x \\ t_y \\ t_z
\end{pmatrix} = \begin{pmatrix}
10 \\ 10 \\ -50
\end{pmatrix}.$$
%
Assume that the Lighthouse base station is located at the origin in world coordinates. We also know that the four photodiodes, as mounted on the VRduino, are located at the following positions (specified in mm):
%
$$p_0 = \begin{pmatrix}
-42 \\ 25 \\ 0
\end{pmatrix}\ \
p_1 = \begin{pmatrix}
42 \\ 25 \\ 0
\end{pmatrix}\ \
p_2 = \begin{pmatrix}
42 \\ -25 \\ 0
\end{pmatrix}\ \
p_3 = \begin{pmatrix}
-42 \\ -25 \\ 0
\end{pmatrix}.$$
%
As the Lighthouse base station sweeps its horizontal and vertical laser lines through the room, each of the four photodiodes will be triggered at a particular time stamp, measured in clock ticks of the microcontroller. Given the parameters above, what are these time stamps for all four photodiodes? List both horizontal and vertical time stamps for each photodiode, i.e., 8 values in total.

For your solutions, round the clock ticks to the nearest integer. Assume that the microcontroller runs at 48~MHz. Similar to the lecture and course notes, we define the sequence of rotations from local to world coordinates as yaw-pitch-roll, i.e., $R= R_z \left( \theta_z \right) R_x \left( \theta_x \right) R_y \left( \theta_y \right)$.Note: You may use your preferred tool for this problem (e.g., Matlab, Python, etc.). There is no need to submit code, but do show your derivations in your write-up.

%%%%%%%%%%%%%%%%%%%%%%%%%%%%%%%%%%%%%%%%%%%%%%%%%%%%%%%%%%%%%%%%%%%%%%%%%%%%%%%%%%%%%%%%%%%%%%%%%%%%%%%%%%%%%%%%%%%%%%%%
% Answer

\begin{answer}
\rule{\textwidth}{0.4pt}
\textbf{Answer:}

Write your answer to this question here.

\end{answer}

%%%%%%%%%%%%%%%%%%%%%%%%%%%%%%%%%%%%%%%%%%%%%%%%%%%%%%%%%%%%%%%%%%%%%%%%%%%%%%%%%%%%%%%%%%%%%%%%%%%%%%%%%%%%%%%%%%%%%%%%
%%%%%%%%%%%%%%%%%%%%%%%%%%%%%%%%%%%%%%%%%%%%%%%%%%%%%%%%%%%%%%%%%%%%%%%%%%%%%%%%%%%%%%%%%%%%%%%%%%%%%%%%%%%%%%%%%%%%%%%%
\newpage
\subsection{Robustness of the Inverse Method \points{15}}

Usually, we do not know the ground truth pose of a tracked object, such as the VRduino. To compute it from the measured time stamps, we can use the homography method as discussed in class and in the course notes. In this case, we need to construct a linear system of equations $\boldsymbol{b}=\boldsymbol{A} \boldsymbol{h}$ and solve it for $\boldsymbol{h}$. This requires inverting the matrix $\boldsymbol{A}$. As discussed in class, the robustness of a solution to such a linear inverse problem (with respect to sensor noise or slight errors in the measurements) is defined by the condition number of the matrix. As before, can may your preferred tool for the following computations (e.g., Matlab, Python, etc.). There is no need to submit code, but show your derivations in your write-up.

\begin{enumerate}[label=(\roman*)]
\item Show the matrix $\boldsymbol{A}$ constructed from the eight measurements you calculated in Section~\ref{sec:imageformation}. \pointsin{5}
\item Compute the singular values $\sigma_{1,\ldots,8} \left( \boldsymbol{A} \right)$ and the condition number $\kappa \left( \boldsymbol{A} \right) = \frac{ \sigma_{max} \left( \boldsymbol{A} \right) }{ \sigma_{min} \left( \boldsymbol{A} \right) }$ for this matrix. Briefly discuss what this number means for the robustness of your solution to this inverse problem (with respect to noise and measurement errors). \pointsin{10}
\end{enumerate}

%%%%%%%%%%%%%%%%%%%%%%%%%%%%%%%%%%%%%%%%%%%%%%%%%%%%%%%%%%%%%%%%%%%%%%%%%%%%%%%%%%%%%%%%%%%%%%%%%%%%%%%%%%%%%%%%%%%%%%%%

\begin{answer}
\rule{\textwidth}{0.4pt}
\textbf{Answer:}

\begin{enumerate}[label=(\roman*)]
\item Write your answer to this question here.

\item Write your answer to this question here.
\end{enumerate}
\end{answer}



%%%%%%%%%%%%%%%%%%%%%%%%%%%%%%%%%%%%%%%%%%%%%%%%%%%%%%%%%%%%%%%%%%%%%%%%%%%%%%%%%%%%%%%%%%%%%%%%%%%%%%%%%%%%%%%%%%%%%%%%
%%%%%%%%%%%%%%%%%%%%%%%%%%%%%%%%%%%%%%%%%%%%%%%%%%%%%%%%%%%%%%%%%%%%%%%%%%%%%%%%%%%%%%%%%%%%%%%%%%%%%%%%%%%%%%%%%%%%%%%%
\newpage
\subsection{Arranging Photodiodes in 3D \points{20}}

On the VRduino, all photodiodes are arranged in the plane of the device. While this makes the pose calculations a bit easier, it may not be the best approach when robustness and precision of the tracking matter. Indeed, the arrangement of the photodiodes on the controllers and the HMDs that Lighthouse is designed to track often have much more complex 3D configurations. Let's look at such a 3D photodiode arrangement in more detail. Once again, you can use your favorite tool for the following computations (e.g., Matlab, Python, etc.). There is no need to submit code, but show your derivations in your write-up.

\begin{enumerate}[label=(\roman*)]
\item What is the minimum number of photodiodes that must be arranged in a non-planar 3D configuration to result in a ``square'' or ``tall'' matrix $\boldsymbol{A} \in \mathbb{R}^{m \times n}$, for $m \geq n$, in the linear system of equations? How did you come up with that number? \pointsin{10}
\item Assume that we have the same four co-planar photodiodes listed in Section~\ref{sec:imageformation} and two additional photodiodes that extrude from the device plane and are located at:
%
$$p_4 = \begin{pmatrix}
0 \\ 25 \\ 10
\end{pmatrix}\ \
p_5 = \begin{pmatrix}
0 \\ -25 \\ -10
\end{pmatrix}.$$
%
On this modified VRduino, assume we have measured the following horizontal and vertical time stamps for all 6 photodiodes and have converted them into normalized coordinates:
%
\begin{center}
\begin{tabular}{ c|c|c }
\hline \hline
photodiodes & $x^n$ & $y^n$ \\
 \hline
 $p_0$ & -0.2926 & -0.0822 \\
 $p_1$ & 0.3296 & 0.9955 \\
 $p_2$ & 0.6459 &  0.4924 \\
 $p_3$ & 0.1919 & -0.2940 \\
 $p_4$ & 0.0948 & 0.4814 \\
 $p_5$ & 0.3403 & 0.1154 \\
 \hline \hline
\end{tabular}
\end{center}
%
Using the least squares solution $\boldsymbol{h} = (\boldsymbol{A}^T \boldsymbol{A})^{-1} \boldsymbol{A}^T \boldsymbol{b}$ to the linear system, retrieve the translation vector $\vec{t}$ and report the yaw, pitch, and roll angles in degrees. Round translations and angles (in degrees) to the nearest integer. \pointsin{10}

\textbf{Hint:} First derive the linear system, invert it to obtain the rotation matrix $R$ and translation vector $\vec{t}$ based on $\boldsymbol{h}$. Then refer to Appendix~B of the course notes for details on how to calculate yaw, pitch, and roll from $R$.
\end{enumerate}

%%%%%%%%%%%%%%%%%%%%%%%%%%%%%%%%%%%%%%%%%%%%%%%%%%%%%%%%%%%%%%%%%%%%%%%%%%%%%%%%%%%%%%%%%%%%%%%%%%%%%%%%%%%%%%%%%%%%%%%%

\begin{answer}
\rule{\textwidth}{0.4pt}
\textbf{Answer:}

\begin{enumerate}[label=(\roman*)]
\item Write your answer to this question here.

\item Write your answer to this question here.
\end{enumerate}
\end{answer}


%%%%%%%%%%%%%%%%%%%%%%%%%%%%%%%%%%%%%%%
%  Programming Written Deliverables  %
%%%%%%%%%%%%%%%%%%%%%%%%%%%%%%%%%%%%%%%
\newpage
\section*{Programming Part PDF Deliverables}
\subsubsection*{2.1.2 Assessing Your Marker Tag}
What are the limitations you observe with your marker-based tracking? How might these limitations be addressed, either with improved software or hardware components? Do you think marker-based tracking could be adopted for commercial systems? If so, why don't you think it's commonly used in products? If not, what limitations make marker-based tracking a poor choice and how do other systems address these limitations?

\begin{answer}
\rule{\textwidth}{0.4pt}
	
	\textbf{Answer:}
	
	Write your answer to this question here.

\rule{\textwidth}{0.4pt}
\end{answer}

\subsubsection*{2.2.2 Tuning the Positional Tracker Filter}
Report the positional filter value that you prefer (i.e., the value that suppresses positional tracking errors, without introducing objectionable latency).

\begin{answer}
\rule{\textwidth}{0.4pt}
	
	\textbf{Answer:}
	
	Write your answer to this question here.

\rule{\textwidth}{0.4pt}
\end{answer}

\subsubsection*{2.2.3 Assessing your Positional Tracking Results}
Comment on the quality of marker-based orientation tracking. What could we change about the camera and/or markers to address the limitations you observe?

comment on the quality of your positional tracking. How does it compare to commercial systems you're tried? What are the primary issues you observe, even after tuning the positional and orientation tracking filters? How might you improve the performance of 6-DOF tracking with both software and hardware upgrades to the CSE 480V kit, assuming it is still based on marker tags for positional tracking?

\begin{answer}
\rule{\textwidth}{0.4pt}
	
	\textbf{Answer:}
	
	Write your answer to this question here.

\rule{\textwidth}{0.4pt}
\end{answer}

\subsubsection*{2.3.1 Calibrating Your Camera using OpenCV}
Include your calibration parameters and details about your webcam (e.g., 2019 13-inch MacBook Pro).

\begin{answer}
\rule{\textwidth}{0.4pt}
	
	\textbf{Answer:}
	
	Write your answer to this question here.

\rule{\textwidth}{0.4pt}
\end{answer}

\subsubsection*{2.3.2 Assessing your Camera Calibration Results} 
Explain why the 9 values of the recovered \texttt{camera\_matrix} in \texttt{camera.yml} are reasonable. Hint: Estimate the field of view for your webcam and relate this to the values in the camera matrix. What else can you conclude about the construction of your webcam from the calibration results?

\begin{answer}
\rule{\textwidth}{0.4pt}
	
	\textbf{Answer:}
	
	Write your answer to this question here.

\rule{\textwidth}{0.4pt}
\end{answer}

\end{document}


\end{document}